\begin{qcard}{2099--2100 \textbullet\ Q0(a)}{Course \textbullet\ TF page}{10}
A question card. The coral badge on the right shows the printed marks. Figures from the original paper are redrawn in vector graphics inside the card.
\end{qcard}
\begin{solution}
The solution panel always follows its question: derivations, intermediate tables and short explanations. Key results are highlighted:
\ans{like this.}
\end{solution}
\repeated{2}
\begin{qcard}{Paper A}{Q1}{12}A question asked in a recent session\ldots\end{qcard}
\begin{qcard}{Paper B}{Q4}{15}\ldots and asked again in an earlier session.\end{qcard}
\begin{solution}One shared solution covers every stacked card.\end{solution}
\begin{notebox}Gold boxes flag a misprint, an illegible part of the scanned paper, or an assumption made in the solution.\end{notebox}

\medskip{\sffamily\small\bfseries\color{slate} Conventions}
\begin{itemize}[leftmargin=1.4em,itemsep=2pt]
\item \textbf{Scope:} 2018--19 to 2023--24. Course codes: \textbf{CSE 219} (current), \textbf{EEE 211/301} (the earlier course on the same content; ``EEE p.$N$'' is the page in the EEE paper collection), \textbf{CSE 311} (Data Communication; only the signals/Fourier/sampling parts are included). CSE 219 2023--24 pages are in the ``22 batch'' file.
\item \textbf{Ordering:} within a topic, the most recent session first.
\item \textbf{Notation:} $\F$, $\Lap$, $\Z$ are the Fourier, Laplace and $z$ transforms; $r(t)=t\,u(t)$; $\operatorname{sinc}x=\sin(\pi x)/(\pi x)$ unless a solution says otherwise; $W_N=e^{-j2\pi/N}$.
\item \textbf{Syllabus:} Anik Sir excludes Parseval's theorem and slides marked $(*)$; Parseval questions are still solved here because they were asked, and are marked in a note.
\item \textbf{Figures:} some scanned figures are partly illegible; the reading used is stated in a note before the solution.
\item \textbf{Last chapter:} state-space and mechanical-analogue questions are not on either teacher's slides; they are listed but not solved.
\end{itemize}
